\section{Conclusiones}
~\label{sec:conclusiones}

Los estándares maduros del monitoreo de red, como la RFC 5470 sobre IPFIX, 
fueron diseñados para entornos de telecomunicaciones de gran escala, donde 
los \textit{observation points} son routers empresariales y los 
\textit{collectors} son piezas dedicadas de infraestructura. Aplicarlos a 
una red LAN de decenas de endpoints monitoreada por un servidor doméstico 
obliga a una traducción que los estándares no documentan: qué componentes 
se conservan, cuáles se simplifican, cuáles pierden sentido. El sistema 
demuestra que esa traducción es posible, pero también que la literatura 
actual no cubre los despliegues intermedios entre el laboratorio didáctico 
y la infraestructura de operador. Caracterizar con rigor ese vacío podría 
ser, por sí solo, el objeto de un trabajo posterior.

El marco legal colombiano aplicable al monitoreo de tráfico en redes 
privadas, compuesto por la Ley 1273 de 2009 sobre delitos informáticos, 
la Ley 1581 de 2012 sobre protección de datos personales, el Decreto 1377 
de 2013 que la reglamenta, y la Ley 2466 de 2025 que actualiza el deber 
de aviso de privacidad, ofrece un perímetro amplio pero deliberadamente 
genérico. El proyecto asumió la responsabilidad de traducir ese perímetro 
a decisiones técnicas concretas que incluyen sanitización del payload 
para mantenerse dentro de la finalidad declarada de monitoreo, recorte 
del contenido sensible para satisfacer el principio de minimización, 
consentimiento del usuario en el endpoint mediante un archivo explícito 
de aceptación, y delimitación del uso a entornos controlados con fines 
académicos. La conclusión que se impone no es que el sistema cumpla la 
norma, sino que el cumplimiento efectivo requiere esa traducción en donde 
el marco legal por sí solo no indica cuántos bytes debe conservar un 
sensor ni cómo debe construirse un flujo de consentimiento. Ese trabajo 
intermedio, que en este caso se realizó dentro del proyecto, no es 
habitual en la literatura técnica sobre monitoreo y queda como 
contribución documentada del informe. La totalidad de las actividades 
descritas se desarrolló dentro de un entorno controlado montado para los 
fines de este proyecto integrador, y la información recopilada durante 
las pruebas tuvo carácter estrictamente investigativo, sin transferencia 
a terceros, sin integración a entornos productivos y sin uso alguno ajeno 
a la validación académica del sistema.

El cronograma de una metodología por fases determina cuándo cada decisión 
debe estar documentada, pero la evidencia técnica que valida o 
descarta esa decisión sigue un ritmo propio que no siempre 
coincide con los hitos; una decisión de arquitectura consignada 
en Elaboración puede quedar invalidada por una prueba ejecutada 
semanas después en Construcción, y el artefacto que la documenta 
ya no describe el sistema real. De ello, se deduce que una
 metodología útil en este contexto no es la que prescribe la 
 secuencia más completa sino la que admite que los artefactos 
 evolucionen de forma versionada y cíclica, y que las decisiones 
 inicialmente tentativas queden registradas como tales hasta que 
 la evidencia técnica las confirme o las revise.

La expresión \textit{tiempo real} sugiere inmediatez, pero en un sistema 
de monitoreo siempre existe un retardo compuesto por varias etapas que 
procesan los datos antes de mostrarlos. Reducir ese retardo en cualquiera 
de esas etapas tiene una contrapartida técnica, que puede manifestarse 
como mayor volumen de falsas alarmas, un mayor consumo de recursos o 
menor calidad del análisis. La decisión no está en si el sistema alcanza 
o no el tiempo real, sino con qué consciencia se interpreta cada alerta. 
Los sistemas de monitoreo no fallan cuando el retardo existe, fallan 
cuando el retardo se ignora y se toman decisiones operativas como si la 
información recibida describiera el presente exacto de la red.

El valor del ejercicio no radicó en producir una innovación técnica 
aislada, sino en articular decisiones de ingeniería, restricciones 
normativas y elecciones de diseño bajo un marco metodológico trazable. 
Lo que el proyecto aporta es la cohesión con que esas piezas se 
combinaron para un problema concreto, y la transparencia con que las 
decisiones se documentaron para que el proceso sea auditable por 
terceros. Esa transparencia, más que la solución en sí, es lo que 
distingue un proyecto académico e investigativo del desarrollo de un 
producto comercial.